\textbf{Status:} PILOT-LOCKED / UPDATED THROUGH v0.6 REVIEW\\
\textbf{Date:} 21 September 2026

This appendix freezes the cheapest experiments needed before substantive adaptation compute. Items explicitly reopened below must be re-frozen under a new pilot register before result-bearing P3/adaptation runs.

\subsection{P1 --- H-Neuron replication}\label{p1-h-neuron-replication}

\begin{itemize}
\tightlist
\item Model: meta-llama/Llama-3.1-8B-Instruct; exact revision pinned at first run.
\item Identification: TriviaQA, CETT answer-token profiles, L1 sparse logistic regression.
\item Thresholds: 0.0005 / 0.001 / 0.002; primary 0.001.
\item OOD: NQ-Open, BioASQ; NonExist diagnostic.
\item Sample: 300 eligible/set if available.
\item Seeds: 42, 314, 2718 where applicable.
\item G0-A remains a pilot compute-screening gate, not a scientific replication threshold.
\end{itemize}

\subsection{P1b --- Causal phenotype}\label{p1b-causal-phenotype}

\begin{itemize}
\tightlist
\item Primary: FaithEval Counterfactual Context.
\item Greedy decoding, max 256 new tokens, rule-based parser.
\item H alpha=0.0 primary; alpha=0.5 sensitivity.
\item Controls: sham, random, same-layer, contribution-matched.
\item G0-B: >=5 pp compliance reduction vs baseline.
\item G0-C: >=3 pp advantage vs contribution-matched control; 95\% bootstrap CI excludes zero.
\item G0-D: relative PPL increase <10\%.
\item Full sparse-probe solver/normalization/convergence settings must be frozen before result-bearing P1; multi-seed selection stability is reported rather than hidden.
\end{itemize}

\subsection{P2 --- Persistent mask / conditional t0 backup audit}\label{p2-persistent-mask-t0}

\begin{itemize}
\tightlist
\item Zero mask after gated MLP intermediate activation and before down-projection.
\item Audit abs(masked value) <=1e-7.
\item Rerun localization and freeze t0 alternative-locus table before any training.
\item Add marginal and conditional causal effects under the H-mask.
\item Add conditional co-ablation / backup discovery where feasible, using a bounded candidate pool if neuron-scale exhaustive analysis is prohibitive.
\item Add residual-stream decodability as a diagnostic, not a pass/fail gate.
\item Freeze a practical-equivalence margin for B0-A near-null backups before adaptation.
\end{itemize}

\subsection{P3 --- Recoverability}\label{p3-recoverability}

\begin{itemize}
\tightlist
\item Synthetic integer arithmetic, modular arithmetic, Boolean/symbolic logic.
\item Calibrate dev difficulty before probe fitting; report difficulty-stratified results.
\item Natural trajectory collection remains deterministic.
\item I0 primary: exact pre-stop prefix, stochastic continuation, EOS allowed.
\item I1 secondary: forced-compute continuation with EOS constrained.
\item I2 robustness: fixed procedural re-check instruction with no new task facts.
\item Do not pool I0/I1/I2 into one recoverability label.
\item K is REOPENED. Simulate K in {8,16,32} and freeze by a pre-specified precision/cost rule before P3.
\item Primary predictive score: held-out continuation-level Bernoulli / hierarchical-binomial negative log likelihood or deviance.
\item Baselines include prompt-only and question-only hidden-state predictors before J-space.
\end{itemize}

\subsection{Stop rules}\label{stop-rules}

If Gate 0 fails, do not run constrained adaptation as a reconstitution experiment. If the mask audit fails, patch under a new protocol version. If conditional t0 backups already explain recovery, classify self-repair rather than reconstitution. If recoverability does not beat simple baselines, deprioritize J-space causal work.
